\section{预计的研究难点及解决方案}

\begin{enumerate}
  \item \textbf{数据集和标签构造难度。}公开Agent安全数据集中，多数侧重提示注入或最终攻击成功率，未必包含完整的工具调用链、权限标签和风险路径标签。若完全依赖公开数据，可能无法支撑本课题的指标计算和归因实验。拟采用"公开基准思想+自建可控任务集"的方式：参考AgentDojo\cite{balunovic2024agentdojo}、OSWorld\cite{xie2024osworld}等基准的任务组织方式，但使用本地可控工具链构造日志；标签由规则生成后再进行人工抽查，保证可复现性和标注成本可控。

  \item \textbf{指标权重的平衡与校准。}六个风险指标（$S_{\text{tool}}, D_{\text{call}}, P_{\text{diff}}, I_{\text{irrev}}, E_{\text{expo}}, H_{\text{cov}}$）的量纲和动态范围不同，直接加权组合可能导致某些信号主导综合风险评分。拟在已知风险标签的验证集上校准各指标的权重参数，使综合风险分数与真实标签的相关性最大化；同时设计消融实验评估各指标的独立贡献，对低信息量指标进行降权或合并处理。

  \item \textbf{行为轨迹图构建的完整性与噪声。}不同Agent执行框架的日志字段存在缺失或不一致，可能导致构建的轨迹图缺失部分依赖边，影响检测和归因的准确性。拟设计鲁棒的图构建规则：对于缺失字段，采用默认值填充并标注不确定性；对于参数来源不明的情况，增加"unknown\_source"节点，保留图结构的完整性，同时在评估中排除低置信度样例的干扰。

  \item \textbf{归因结果可信度的验证。}风险归因本身具有一定主观性，若没有可验证标签，可能难以证明归因结果正确。拟在构造数据时同步记录真实风险触发节点和风险路径，形成gold path与gold node标签；评价时使用Path Recall@K、Hit@K和案例分析共同验证归因结果，以定量方式证明方法的有效性。

  \item \textbf{论文创新性表达与工程实现的平衡。}单纯做指标汇总或日志可视化可能被认为工程性较强、方法创新不足。拟将研究点二作为会议论文主线，突出行为轨迹图、图模式检测、风险路径定位和反事实归因四项方法贡献；研究点一作为基础框架提供评估基座，通过消融实验和理论分析凸显其在序列层面对组合风险建模的创新价值。
\end{enumerate}
